% ID: FIS-DIN-NEW-001
% Titolo: Due forze perpendicolari e moto nel piano
% Disciplina: Fisica
% Area: Dinamica
% Argomento: Leggi di Newton
% Sottoargomento: Moto piano con accelerazione costante
% Classe: Terza liceo scientifico
% Difficolta: 4
% Tipo: problema_numerico
% Risultato: a = (3,0; -4,0) m/s^2; t = 0,60 s; x = 1,02 m; v = (2,6; -1,2) m/s
% Tempo_stimato: 12 min
% Competenze: scomposizione vettoriale, modellizzazione, cinematica bidimensionale
% Prerequisiti: secondo principio della dinamica, moto uniformemente accelerato, vettori
% Tag: fisica, dinamica, leggi_newton, forze_perpendicolari, moto_piano
% Fonte: verifica_fisica_dinamica_anonimizzata
% Autore: Riccardo Carli
% Licenza: CC-BY-SA-4.0

\begin{esercizio}
Un corpo di massa \(m = 0{,}60\,\mathrm{kg}\) e soggetto a due forze costanti
e perpendicolari tra loro: \(F_1 = 1{,}8\,\mathrm{N}\), diretta lungo
\(+\hat x\), e \(F_2 = 2{,}4\,\mathrm{N}\), diretta lungo \(-\hat y\).

All'istante iniziale il corpo si trova nell'origine e possiede velocita
iniziale
\[
\vec v_0 =
\begin{pmatrix}
0{,}80 \\[4pt]
1{,}20
\end{pmatrix}
\mathrm{m/s}.
\]

Determina l'accelerazione del corpo, l'istante in cui esso attraversa
nuovamente l'asse \(x\), cioe quando \(y=0\) per la seconda volta, e la
coordinata \(x\) e la velocita in quell'istante.
\end{esercizio}

\begin{soluzione}
La risultante delle forze e:
\[
\vec F_\text{tot}=
\begin{pmatrix}
1{,}8\\[4pt]
-2{,}4
\end{pmatrix}
\mathrm{N}.
\]

Dal secondo principio della dinamica:
\[
\vec a=\frac{\vec F_\text{tot}}{m}
=\frac{1}{0{,}60}
\begin{pmatrix}
1{,}8\\[4pt]
-2{,}4
\end{pmatrix}
=
\begin{pmatrix}
3{,}0\\[4pt]
-4{,}0
\end{pmatrix}
\mathrm{m/s^2}.
\]

Le equazioni orarie sono:
\[
x(t)=v_{0x}t+\frac{1}{2}a_xt^2
=0{,}80t+1{,}5t^2,
\]
\[
y(t)=v_{0y}t+\frac{1}{2}a_yt^2
=1{,}20t-2{,}0t^2.
\]

Per trovare il secondo attraversamento dell'asse \(x\), imponiamo \(y=0\):
\[
1{,}20t-2{,}0t^2=0.
\]
\[
t(1{,}20-2{,}0t)=0.
\]
La soluzione \(t=0\) e l'istante iniziale; quella successiva e:
\[
t=\frac{1{,}20}{2{,}0}=0{,}60\,\mathrm{s}.
\]

La coordinata \(x\) in quell'istante e:
\[
x(0{,}60)=0{,}80\cdot 0{,}60+1{,}5\cdot (0{,}60)^2
=1{,}02\,\mathrm{m}.
\]

La velocita e:
\[
\vec v(t)=\vec v_0+\vec a t.
\]
Quindi:
\[
\vec v(0{,}60)=
\begin{pmatrix}
0{,}80\\[4pt]
1{,}20
\end{pmatrix}
+
0{,}60
\begin{pmatrix}
3{,}0\\[4pt]
-4{,}0
\end{pmatrix}
=
\begin{pmatrix}
2{,}6\\[4pt]
-1{,}2
\end{pmatrix}
\mathrm{m/s}.
\]

Il modulo della velocita in quell'istante e:
\[
|\vec v|\simeq 2{,}9\,\mathrm{m/s}.
\]
\end{soluzione}
